\chapter{Literature Review}
\label{chap:literature-review}
\chapterrunning{literature review}

\section{Heuristics, Hyper-Heuristics, and Algorithm Selection}
\label{sec:lit-heuristics}

The practical difficulty of many combinatorial optimization problems has long made heuristic design
an essential part of computer science and operations research. Rice's classic formulation of the
algorithm selection problem cast the issue in its most general form: given a space of instances, a
portfolio of candidate algorithms, and a performance measure, the task is to learn or construct a
mapping from instances to algorithms that performs well in expectation \cite{rice1976}. In modern
notation, if $\mathcal{X}$ is an instance family, $\mathcal{A}$ is a portfolio of algorithms, and
$m(a,x)$ is a lower-is-better performance metric, then the ideal selector is
\begin{equation}
\pi^{\star} \in \argmin_{\pi : \mathcal{X} \to \mathcal{A}}
\mathbb{E}_{x \sim \mathcal{D}}[m(\pi(x),x)].
\label{eq:lit-alg-selection}
\end{equation}
Equation~\eqref{eq:lit-alg-selection} remains foundational because it separates the underlying
optimization problem from the meta-level question of which procedure should be applied to which
instance.

Hyper-heuristics emerged as one prominent response to this meta-level challenge. Rather than
searching directly over candidate solutions, hyper-heuristics search over heuristics or heuristic
components, with the goal of automating heuristic selection, sequencing, or generation
\cite{burke2013,dokeroglu2024}. This shift from solution space to heuristic space is especially
relevant for the present dissertation because the proposed LLM-guided loop also operates primarily
at the level of \emph{programs that produce solutions}, not at the level of solutions themselves.
Recent surveys emphasize that hyper-heuristics now span selection, generation, portfolio, and
configuration settings, and that modern work increasingly integrates machine learning into all four
categories \cite{dokeroglu2024}. The present dissertation belongs to this tradition, with the
difference that the search operator is a language model capable of proposing new executable code
rather than a classical mutation, grammar, or rule-combination mechanism.

Algorithm selection has likewise evolved toward both feature-based and feature-free approaches.
Traditional methods rely on hand-designed instance descriptors, whereas newer methods attempt to
learn directly from raw instance representations \cite{alissa2023}. This distinction becomes
important in later chapters, especially for the adaptive heuristic portfolio experiments, because
it frames a key design choice: should a controller rely on explicit descriptors crafted from domain
knowledge, or can a more generic representation learn the selector implicitly?

\section{Formal Benchmark Problem Definitions}
\label{sec:lit-problems}

The empirical phases of this dissertation are built around standard benchmark optimization
families. A concise formal account is therefore useful before turning to learning-based methods.

\subsection{Symmetric Traveling Salesperson Problem}
\label{sec:lit-tsp}

In the symmetric traveling salesperson problem (TSP), one is given a complete weighted graph
$G = (V,E,w)$ with $w_{ij} = w_{ji} \geq 0$ and seeks a minimum-cost Hamiltonian cycle
\cite{reinelt1991}. Using binary edge-selection variables $x_{ij}$, a standard formulation is
\begin{align}
\min_{x} \quad & \sum_{(i,j)\in E} w_{ij}x_{ij} \label{eq:lit-tsp-obj} \\
\text{s.t.} \quad
& \sum_{j \neq i} x_{ij} = 2, && \forall i \in V, \label{eq:lit-tsp-degree}\\
& \sum_{i,j \in S,\ i<j} x_{ij} \le |S|-1, && \forall S \subset V,\ 2 \le |S| < |V|, \label{eq:lit-tsp-subtour}\\
& x_{ij} \in \{0,1\}, && \forall (i,j)\in E. \label{eq:lit-tsp-binary}
\end{align}
Constraint \eqref{eq:lit-tsp-degree} enforces degree two at every city, while
\eqref{eq:lit-tsp-subtour} removes disconnected subtours. The TSP is NP-hard, which explains why
heuristic methods remain central despite decades of exact and approximate progress.

\subsection{Asymmetric Traveling Salesperson Problem}
\label{sec:lit-atsp}

The asymmetric TSP (ATSP) replaces the undirected graph by a complete directed graph
$G = (V,A,c)$ in which $c_{ij}$ and $c_{ji}$ may differ. The goal is again to find a minimum-cost
Hamiltonian tour, but now with one incoming and one outgoing arc per city. Many constructive ideas
transfer from TSP to ATSP, yet the asymmetry changes the local-search geometry and often weakens
heuristics that rely on symmetric exchange structure. This makes ATSP a useful transfer family for
testing whether a learned or evolved heuristic reflects a genuine search principle rather than a
benchmark-specific shortcut.

\subsection{Capacitated Vehicle Routing Problem}
\label{sec:lit-cvrp}

In the capacitated vehicle routing problem (CVRP), a depot node $0$, customer set
$V = \{1,\dots,n\}$, customer demands $q_i > 0$, and vehicle capacity $Q$ are given. The task is
to partition customers into depot-rooted routes whose total demand does not exceed $Q$ and whose
combined travel cost is minimal \cite{clarke1964,uchoa2017}. A compact route-based view is
\begin{align}
\min_{\mathcal{R}} \quad & \sum_{R \in \mathcal{R}} c(R) \label{eq:lit-cvrp-obj}\\
\text{s.t.} \quad
& \bigcup_{R \in \mathcal{R}} R = V, \qquad
R \cap R' = \varnothing \text{ for } R \neq R', \label{eq:lit-cvrp-partition}\\
& \sum_{i \in R} q_i \le Q, \qquad \forall R \in \mathcal{R}. \label{eq:lit-cvrp-capacity}
\end{align}
Unlike TSP, CVRP combines sequencing and packing structure. That added constraint pressure matters
for this dissertation because a program that looks promising under tour construction alone may fail
once feasibility management, route merging, and penalty handling become part of the evaluator.

\section{Machine Learning for Combinatorial Optimization}
\label{sec:lit-mlco}

Machine learning for combinatorial optimization has developed into a substantial field in its own
right. Bengio, Lodi, and Prouvost provide one of the clearest organizing views: optimization
problems can be treated as data points, and learning can be used to replace costly or poorly
defined decisions inside larger solving procedures \cite{bengio2021mlco}. This framing is useful
because it avoids the misleading idea that learning must replace classical optimization wholesale.
Instead, machine learning can support branching, ranking, repair, construction, selection, or
parameterization within a broader pipeline.

A complementary research stream studies direct learning-based solvers. Neural combinatorial
optimization with reinforcement learning showed early promise on Euclidean TSP by training neural
policies to construct tours directly from coordinates \cite{bello2016nco}. Khalil et al.\ extended
this direction by learning greedy algorithms over graph-structured state representations, showing
that learned policies can act as reusable meta-algorithms rather than as fixed instance-specific
predictors \cite{khalil2017learning}. Subsequent work extended learning-based optimization to more
constrained settings and more realistic objectives, but the broader lesson remained the same:
learned solvers can be effective, yet strong classical heuristics often remain difficult to beat
across wide instance distributions \cite{bengio2021mlco}. This tension matters for the current
dissertation, which does not ask an LLM to output one final tour directly. Instead, it asks
whether an LLM can iteratively improve the \emph{code} of a heuristic solver under benchmark
feedback.

The routing literature also provides the baselines and datasets that make such claims meaningful.
TSPLIB remains the canonical reference library for traveling salesperson problems
\cite{reinelt1991}. For symmetric TSP, the Lin--Kernighan family and its refinements, including
Helsgaun's implementation, remain among the strongest and most influential heuristic baselines
\cite{lin1973,helsgaun2000}. For CVRP, Clarke and Wright's savings heuristic remains a foundational
constructive baseline \cite{clarke1964}, while modern state-of-the-art practice includes far
stronger hybrid metaheuristics such as Vidal's hybrid genetic search \cite{vidal2022hgs}. The
modern Uchoa instance family likewise provides a more challenging and statistically informative
benchmark regime than many older CVRP sets \cite{uchoa2017}. Together, these references define the
level of baseline strength against which LLM-evolved heuristics must be judged. A language-model
system that only beats naive constructive baselines should not be conflated with one that beats
the best established heuristics in the literature.

\section{Large Language Models for Code Generation}
\label{sec:lit-codegen}

The empirical framework studied in this dissertation depends directly on the rapid improvement of
code-capable language models.
Codex demonstrated that large language models trained on code could solve a meaningful fraction of
functional programming tasks and that repeated sampling substantially boosts success rates on
harder problems \cite{chen2021codex}. Austin et al.\ likewise showed that program synthesis
performance scales strongly with model size and that natural-language feedback can materially
improve generated programs \cite{austin2021program}. These findings indicate that outer-loop search
and repair are integral components of strong performance on demanding code-generation tasks.

As the field matured, surveys began to distinguish between simple code completion, benchmarked
program synthesis, repair, debugging, and agentic multi-step workflows. Early overviews of
NL2Code models emphasized training data, scaling, and benchmark design \cite{zan2023nl2code},
while more recent surveys document a much larger and more heterogeneous landscape of code models,
evaluation protocols, and practical risks \cite{jiang2024codesurvey}. At the same time, stronger
benchmarks such as AlphaCode-style competitive programming and LiveCodeBench reveal that the
difference between solving simple textbook-style tasks and solving harder, more recent coding
problems remains substantial \cite{li2022alphacode,jain2024livecodebench}. This distinction
reinforces the need to evaluate code generation under task-specific execution and held-out tests,
rather than by surface plausibility alone.

\section{Evaluator-Guided Search and Algorithm Discovery}
\label{sec:lit-evaluator-search}

One-shot code generation is only one end of the design spectrum. A more relevant line of work for
this dissertation combines generative models with evaluators that score candidate programs and feed
the results back into an outer search loop. In abstract form, the objective is no longer to
predict source code that merely ``looks right,'' but to search a program space $\mathcal{H}$ for
\begin{equation}
h^{\star} \in \argmin_{h \in \mathcal{H}} J_{\mathcal{D}}(h),
\label{eq:lit-program-search}
\end{equation}
where $J_{\mathcal{D}}$ is defined by measurable performance on a benchmark panel. This framing is
important because it aligns the generative model with verification rather than with free-form
textual plausibility.

FunSearch is a central example. It pairs a frozen LLM with a systematic evaluator and a diverse
program archive in order to discover high-scoring programs for problems that are hard to solve but
easy to evaluate \cite{romeraparedes2024funsearch}. Its importance lies not only in its results,
but in its formulation: the LLM is used as a variation operator inside an evolutionary process,
while the evaluator guards against incorrect or hallucinated proposals.

AlphaDev demonstrates a related idea from a reinforcement-learning perspective. It formulates
algorithm discovery for low-level sorting routines as a single-player game and optimizes directly
for correctness and latency at the assembly level \cite{mankowitz2023alphadev}. AlphaCode, while
not an evolutionary system in the same sense, also depends heavily on large-scale sampling,
behaviour-based filtering, and selection among many candidate programs \cite{li2022alphacode}.
These systems show that modern code-generating AI often relies on the interaction between
generative diversity and external verification, not on a single perfect sample.

The most direct predecessor to the present dissertation in combinatorial optimization is the study
of LLMs as evolutionary optimizers by Liu et al.\ \cite{liu2023lmea}. Their system treats the LLM
as a crossover-and-mutation operator within a population-based optimizer and shows promising TSP
results on very small instances. This work is conceptually important, but its scale also exposes a
gap. There remains a need for research that studies executable whole-program heuristic evolution on
standard benchmark families, with stronger baselines, repeated trials, transfer evaluation, and
mechanism controls. A recent systematic review of LLMs in combinatorial optimization confirms both
the rapid expansion of the area and the diversity of roles LLMs now play in optimization systems
\cite{daros2025llmco}. However, the review also implies that careful cross-family empirical
evaluation remains comparatively rare.

\section{Curriculum, Replay, and Diversity Mechanisms}
\label{sec:lit-replay}

The replay-aware components of this dissertation draw from several adjacent literatures rather than
from one single source. Curriculum learning offers one motivation. Bengio et al.\ formalized the
idea that ordering examples from easier to harder can improve learning and generalization
\cite{bengio2009curriculum}. Later surveys and automated curriculum methods extended the idea to a
wide range of machine learning and reinforcement learning settings
\cite{soviany2021curriculumsurvey,graves2017automatedcurriculum}. In adversarial or sequential
domains, curricula matter because the learner's exposure distribution is part of the effective
training algorithm.

Replay mechanisms supply a second motivation. In reinforcement learning, prioritized experience
replay showed that non-uniform reuse of past transitions can improve learning efficiency relative
to uniform replay \cite{schaul2015prioritized}. In one common formulation, transition $i$ is drawn
with probability
\begin{equation}
P(i) = \frac{p_i^{\alpha}}{\sum_j p_j^{\alpha}},
\label{eq:lit-per}
\end{equation}
where $p_i$ is a priority score and $\alpha$ controls the strength of prioritization.
Later work generalized this idea to sequence-level replay and more structured replay decisions
\cite{brittain2019pser}. Although the present dissertation does not replay transitions in the
standard RL sense, it reuses \emph{instances}, failure cases, and archive summaries as a memory
mechanism for future program proposals. That analogy is strong enough to be informative, but weak
enough that it must be tested rather than assumed. A replay strategy that helps temporal-difference
learning does not automatically imply a successful replay strategy for program-level heuristic
evolution.

Diversity-based search provides a third motivation. Quality-diversity methods such as MAP-Elites
and novelty-search-with-local-competition aim not merely to optimize a single objective, but to
retain a collection of high-performing, behaviourally distinct artifacts
\cite{mouret2015mapelites,pugh2016qd}. This family of ideas is relevant because a code-evolution
loop can easily collapse to narrow local search unless the archive preserves multiple useful
regions of the behavioural space. The present dissertation therefore draws on QD-style intuitions
when it tracks novelty, maintains archives, and evaluates whether diversity corresponds to genuine
adaptation rather than lexical churn.

Finally, self-play and adversarial practice provide a motivating example for why curriculum and
iterative pressure can induce capability growth at all. AlphaGo Zero showed that high-level
performance can emerge from repeated self-play without human demonstrations \cite{silver2017alphagozero}.
Its scale and claims differ substantially from those considered here. Nevertheless, the underlying
intuition remains relevant: iterative interaction can sometimes unlock behaviour that is hard to
obtain from static supervision alone.

\section{Methodological Discipline and the Need for Conservative Claims}
\label{sec:lit-methodology}

An important part of the literature relevant to this dissertation concerns \emph{how} claims are
validated, not only \emph{what} systems are built. Henderson et al.\ document the reproducibility
problems that arise when stochastic learning systems are evaluated with too few seeds, weak
statistical reporting, or poorly matched baselines \cite{henderson2018drltm}. Agarwal et al.\
extend this point by arguing for robust statistical summaries and uncertainty-aware comparisons in
deep RL benchmarking \cite{agarwal2021precipice}. Although these works target reinforcement
learning rather than code evolution, the methodological lesson transfers directly. When a
generative system can produce very different candidate programs across runs, best-case anecdotes are
not enough.

This methodological literature motivates the use of paired offsets, aggregate reports, bootstrap
comparisons, held-out panels, and mechanism-specific controls in the present study. A simple paired
comparison between two conditions $A$ and $B$ can be written as
\begin{equation}
\hat{\Delta}_{A,B} = \frac{1}{n}\sum_{i=1}^{n}\left(m_i^{(A)} - m_i^{(B)}\right),
\label{eq:lit-paired-delta}
\end{equation}
with uncertainty estimated by resampling or other repeated-evaluation procedures. A negative result
obtained under proper controls can be more informative than a seemingly positive result that
confounds memory mechanisms with extra candidate budget.

\section{Synthesis and Research Gap}
\label{sec:lit-gap}

Table~\ref{tab:lit-positioning} summarizes the main research traditions from which this
dissertation draws.

\begin{table}[t]
\centering
\scriptsize
\setlength{\tabcolsep}{4pt}
\begin{tabular}{
  >{\raggedright\arraybackslash}p{1.1in}
  >{\raggedright\arraybackslash}p{0.8in}
  >{\raggedright\arraybackslash}p{0.95in}
  >{\raggedright\arraybackslash}p{0.95in}
  >{\raggedright\arraybackslash}p{1.0in}
}
\toprule
Tradition & Search object & Main strength & Main limitation & Representative refs. \\
\midrule
Classical heuristics and metaheuristics & solutions or route moves & strong benchmarked performance & expensive domain expertise & \cite{lin1973,helsgaun2000,clarke1964,vidal2022hgs} \\
Hyper-heuristics and algorithm selection & heuristic choices or generators & explicit heuristic-level reasoning & relies on crafted operators or features & \cite{rice1976,burke2013,dokeroglu2024} \\
Learning-based combinatorial optimization & policies, rankings, or solver components & learns recurring structure from data & difficult to beat strong classical solvers broadly & \cite{bengio2021mlco,bello2016nco,khalil2017learning} \\
Evaluator-guided program search with LLMs & executable programs & open-ended code generation plus verification & noisy search; mechanism claims need strong controls & \cite{romeraparedes2024funsearch,mankowitz2023alphadev,liu2023lmea} \\
\bottomrule
\end{tabular}
\caption[Dissertation positioning relative to adjacent literatures]{Positioning of the present dissertation relative to adjacent literatures. The thesis sits
closest to evaluator-guided program search, but its experimental logic also depends on benchmarked
routing, hyper-heuristic reasoning, and conservative empirical methodology.}
\label{tab:lit-positioning}
\end{table}

\begin{figure}[t]
\centering
\begin{tikzpicture}[
  box/.style={draw, rounded corners, align=center, text width=1.55in, minimum height=0.74in},
  core/.style={draw, rounded corners, align=center, fill=black!6, text width=1.9in, minimum height=0.9in},
  line/.style={-Latex, thick}
]
\node[core] (core) at (0,0) {Cross-family LLM-guided\\code evolution\\under strict controls};
\node[box] (heur) at (-4.2,1.9) {Classical heuristics\\and benchmark\\baselines};
\node[box] (hyper) at (0,3.0) {Hyper-heuristics\\and algorithm\\selection};
\node[box] (mlco) at (4.2,1.9) {Learning for\\combinatorial\\optimization};
\node[box] (llm) at (-3.6,-2.2) {LLM code generation\\and evaluator-guided\\program search};
\node[box] (replay) at (3.6,-2.2) {Curriculum, replay,\\and quality-diversity\\mechanisms};
\draw[line] (heur) -- (core);
\draw[line] (hyper) -- (core);
\draw[line] (mlco) -- (core);
\draw[line] (llm) -- (core);
\draw[line] (replay) -- (core);
\end{tikzpicture}
\caption[Conceptual synthesis of research threads]{Conceptual synthesis of the research threads integrated by the dissertation. The gap is
not any single ingredient in isolation, but the absence of a cross-family, validator-aware account
of how these ingredients interact when the object of search is executable heuristic code.}
\label{fig:lit-synthesis-map}
\end{figure}

The literature reviewed above supports five observations that matter directly for this thesis.

\begin{enumerate}[leftmargin=2em]
\item Classical optimization already provides strong benchmark libraries and strong heuristic
baselines \cite{reinelt1991,lin1973,helsgaun2000,clarke1964,uchoa2017,vidal2022hgs}.
\item Hyper-heuristics and algorithm selection offer established ways to reason about search in
heuristic space rather than solution space \cite{rice1976,burke2013,dokeroglu2024}.
\item Modern code LLMs are strong enough to generate executable programs, especially when combined
with repeated sampling, testing, and repair
\cite{chen2021codex,austin2021program,li2022alphacode,jain2024livecodebench}.
\item Evaluator-guided algorithm discovery with AI is now credible, but existing success stories
often depend on specialized evaluators, narrow scopes, or problem formulations that differ from
full-program benchmark optimization
\cite{mankowitz2023alphadev,romeraparedes2024funsearch,liu2023lmea}.
\item Curriculum, replay, and diversity mechanisms are theoretically motivated, but their value in
LLM-guided \emph{code evolution} remains an empirical question rather than a settled fact
\cite{bengio2009curriculum,schaul2015prioritized,pugh2016qd}.
\end{enumerate}

The present thesis addresses a gap at the intersection of these points. There is still no clear
empirical account of how LLM-guided code evolution behaves across multiple benchmark families when
evaluated with replicated runs, strong classical baselines, replay alternatives, and strict
validator-based transfer tests. The dissertation therefore centers on a conservative cross-family
study of executable heuristics rather than on isolated positive examples.

\section{Chapter Summary}
\label{sec:lit-summary}

This chapter provided the technical context for the dissertation in four steps. It positioned the
work within algorithm selection and hyper-heuristics, where the central object of search is a
heuristic procedure rather than a single candidate solution. It reviewed the benchmark optimization
families that make the empirical claims meaningful, including the standard TSP, ATSP, and CVRP
formulations and their mature baseline literatures. It also situated the project within modern
code-generation and evaluator-guided program-search research, where LLMs are most credible when
paired with strong external verification. Finally, it showed why replay, curriculum, and diversity
mechanisms must be evaluated under careful experimental discipline rather than inferred from
isolated positive cases. The next chapter turns from that context to the experimental design used
to study the gap identified here.
